% !TeX program = lualatex
\documentclass[10pt,a4paper,twocolumn]{article}

% --- 言語切り替えフラグ (\jatrue: 日本語版, \jafalse: 英語版) ---
% --- 言語切り替えフラグ (外部指定がなければデフォルト日本語) ---
\ifdefined\ifja\else
  \newif\ifja
  \jatrue
\fi

\usepackage{luatexja}
\usepackage{luatexja-fontspec}
\setmainjfont{HaranoAjiMincho}
\setsansjfont{HaranoAjiGothic}

\usepackage{geometry}
\geometry{top=20mm,bottom=20mm,left=16mm,right=16mm,columnsep=6mm}

\usepackage{amsmath,amssymb}
\usepackage{booktabs}
\usepackage{listings}
\usepackage{xcolor}
\usepackage{hyperref}
\usepackage{titlesec}

\hypersetup{
    colorlinks=true,
    linkcolor=black,
    citecolor=black,
    urlcolor=blue
}

\lstset{
    basicstyle=\ttfamily\scriptsize,
    backgroundcolor=\color{gray!10},
    frame=single,
    framerule=0.2pt,
    breaklines=true,
    captionpos=b,
    keepspaces=true,
    columns=flexible
}

\title{\textbf{\Large 
\ifja
多層抽象化の欺瞞：\\なぜソフトウェアセキュリティは枯れた技術とVoid Architecture\texttrademark への回帰を要求するのか
\else
The Fallacy of Layered Abstractions:\\Why Software Security Mandates a Return to Proven Primitives and the Void Architecture\texttrademark
\fi
}}

\author{
    \textbf{xsigil} \\
    Section-9.sh Infrastructure Research Division \\
    \texttt{xsigil@section-9.sh}
}

\date{October 8, 2026}

\begin{document}

\maketitle

\begin{abstract}
\ifja
現代のソフトウェア工学は、開発スループットの向上と引き換えに、過度な多層抽象化（Layered Abstraction）を無批判に受け入れてきた。しかし、抽象化レイヤーの積層はシステム内部の実行セマンティクスを不可視化し、境界間のインピーダンスミスマッチ、巨大なサプライチェーン攻撃面、プロトコル解釈の乖離を構造的にもたらす。本稿では、境界防御（WAF、動的解析スキャナ）に依存するブラックボックスセキュリティの破綻を数学的に証明し、システムを検証可能な「ホワイトボックス」へと回帰させる唯一の解として、数十年研ぎ澄まされた枯れた技術（Unix哲学、カーネルプリミティブ、決定論的ステートマシン）のみで構築する\textbf{Void Architecture\texttrademark}を提唱する。さらに、重厚なコンテナランタイムを排除し、Linuxカーネルの \texttt{unshare(2)}、\texttt{chroot(2)}、\texttt{tmpfs} を直接行使することで、攻撃面を数学的極小値へと漸近させる実証的パラダイムを示す。
\else
Modern software engineering has unreservedly embraced excessive layered abstractions in pursuit of developer velocity. However, this stacking of opaque abstraction layers inevitably obscures execution semantics, leading to impedance mismatches across boundaries, vast supply-chain attack surfaces, and protocol drift. This paper demonstrates the systemic failure of black-box security paradigms reliant on perimeter defenses (WAFs, dynamic scanners) and proposes \textbf{Void Architecture\texttrademark}—constructed solely upon decades-proven primitives (Unix philosophy, direct kernel primitives, and deterministic state machines)—as the sole viable path to verifiable white-box systems. Furthermore, by dismantling heavyweight container engines in favor of direct execution of Linux kernel primitives (\texttt{unshare(2)}, \texttt{chroot(2)}, and \texttt{tmpfs}), we establish an empirical paradigm that asymptotically reduces the attack surface to its mathematical minimum.
\fi
\end{abstract}

\section{
\ifja はじめに \else Introduction \fi
}
\ifja
ソフトウェア産業は過去数十年、Webフレームワーク、高位ORM、サービスメッシュ、コンテナオーケストレータといった抽象化レイヤーの恩恵を謳歌してきた。しかし、この「生産性の向上」は、背後にある物理的実行モデル（システムコール、メモリ配置、ソケットバッファ、バイトストリーム）の不可視化という巨大な代償を支払っている。

抽象化レイヤーが厚くなるほど、開発者は自らのコードが実際にどのような機械語やシステムコール列に変換され、どの境界でどのようにパースされているかを把握できなくなる。この「不可知性（Opaqueness）」こそが、現代のセキュリティ脆弱性が根絶されない本質的な病理である。

本稿では、枯れた技術（Proven Primitives）以外の非本質的抽象化を徹底的に解体し、真のホワイトボックス開発を実現するための設計思想である \textbf{Void Architecture\texttrademark} を論理的・構造的に定式化する。
\else
For decades, the software industry has celebrated the compounding convenience of frameworks, ORMs, service meshes, and container orchestrators. Yet, this surge in developer velocity has extracted a devastating price: the total obfuscation of underlying physical execution models—system calls, memory layouts, socket buffers, and raw byte streams.

As abstraction layers accumulate, developers lose the ability to reason about the exact instructions emitted and how data is parsed across architectural boundaries. This structural opaqueness represents the foundational pathology underlying modern, recurring security vulnerabilities.

This paper formally deconstructs non-essential abstractions and establishes \textbf{Void Architecture\texttrademark}, a design methodology grounded purely in proven primitives to mandate absolute white-box visibility.
\fi

\section{
\ifja 現代的抽象化の構造的破綻 \else The Breakdown of Modern Abstraction \fi
}

\subsection{
\ifja プロトコル解釈の乖離とインピーダンスミスマッチ \else Impedance Mismatches \& Protocol Drift \fi
}
\ifja
現代の分散トポロジにおいて、単一のリクエストはCDN、リバースプロキシ、ロードバランサ、APIゲートウェイ、そして内部ランタイムへとリレーされる。各レイヤーは異なるパースエンジンを実装しており、RFC仕様の微妙な解釈差（例えば \texttt{Content-Length} と \texttt{Transfer-Encoding} の優先順位や多重ヘッダの扱い）が必然的に生じる。

このインピーダンスミスマッチは、HTTP Request Smugglingやキャッシュポイズニングといった致命的な攻撃を構造的に誘発する。レイヤーを多重化する行為そのものが、脆弱性を生み出す境界歪みを不可避に増幅させている。
\else
In typical distributed topologies, a single request traverses front-end CDNs, reverse proxies, load balancers, API gateways, and backend runtimes. Each layer relies on distinct parsing engines, yielding subtle discrepancies in RFC interpretation—such as conflicting handling of \texttt{Content-Length} vs. \texttt{Transfer-Encoding} or duplicate headers.

This protocol drift structurally invites HTTP Request Smuggling and cache poisoning. The compounding of intermediary layers does not enhance defense; it inherently expands the vector space for exploitable boundary desynchronization.
\fi

\subsection{
\ifja ORMによる偽りの安全性 \else The Illusion of Safety via ORM \fi
}
\ifja
オブジェクト関係マッピング（ORM）や抽象クエリビルダは、SQLインジェクションを自明に防御する銀の弾丸として普及した。しかし、ORMが隠蔽したリレーショナル代数のセマンティクスは、より致命的な論理欠陥を胚胎させた。
\begin{enumerate}
    \item \textbf{Mass Assignment}: デシリアライズされた入力ペイロードを内部モデルへ無検証にバインドすることによる、管理者フラグや権限属性の改ざん。
    \item \textbf{Insecure Direct Object Reference (IDOR)}: 動的クエリビルダのメソッドチェーンにおける、マルチテナント境界（\texttt{tenant\_id}）フィルタの意図しない脱落。
\end{enumerate}
生のリレーショナル代数を直接記述することを放棄した結果、プログラマは自らが発行するクエリの決定論性を喪失している。
\else
Object-Relational Mapping (ORM) and dynamic query builders were championed as panaceas against SQL injection. However, by concealing the semantics of relational algebra, ORMs introduced more insidious classes of logical vulnerabilities:
\begin{enumerate}
    \item \textbf{Mass Assignment}: Blind binding of deserialized JSON payloads into database models, enabling unauthorized mutation of privilege flags.
    \item \textbf{Insecure Direct Object Reference (IDOR)}: Accidental truncation of \texttt{tenant\_id} scoping predicates in chained dynamic abstractions.
\end{enumerate}
By surrendering direct control over relational queries, developers forfeit deterministic visibility into database interactions.
\fi

\subsection{
\ifja サプライチェーンの爆発と動的バインディング \else Supply Chain Explosion \& Dynamic Bindings \fi
}
\ifja
一般的なモダンアプリケーションは、直接的・間接的に膨大な外部パッケージに依存している。各依存パッケージの深層には、検証されていない動的C拡張や不可逆なビルド時スクリプトが存在する。ビルドの非決定論性とランタイムの動的リンクは、ソフトウェアサプライチェーン攻撃に対する完全な無防備状態を作り出している。
\else
Contemporary applications routinely import thousands of transitive packages across disparate package registries. Buried deep within these dependency trees reside unvetted native C bindings and opaque build scripts. Build non-determinism coupled with dynamic shared-library linking leaves runtime environments defenceless against supply-chain compromises and remote execution implants.
\fi

\section{
\ifja ブラックボックスセキュリティの破綻 \else The Failure of Black-Box Security \fi
}
\ifja
セキュリティ業界は長年、この抽象化の歪みを「外付けの防御壁」で覆い隠そうとしてきた。しかし、境界防御（WAF、次世代FW）やブラックボックス診断（DAST、外部侵入テスト）には原理的限界が存在する。

\begin{equation}
    P(\text{Breach}) = 1 - \prod_{i=1}^{n} (1 - \epsilon_i)
\end{equation}

ここで、$\epsilon_i$ はレイヤー $i$ におけるパーサの不整合率である。レイヤー数 $n$ が増加するほど、外部WAFが「安全」と判定したパケットが内部ランタイムで悪意ある解釈へと変異する確率 $P(\text{Breach})$ は $1$ に収束する。外部からのパケット検査によって内部状態の論理的整合性を保証することは、計算量理論的にも不可能である。

セキュリティを担保するための唯一の論理的帰結は、\textbf{「内部状態・遷移規則・境界フローが100\%観測および検証可能なホワイトボックス開発」}への回帰である。
\else
For years, the industry has attempted to conceal architectural defects behind perimeter appliances. However, perimeter defenses (WAFs, NGFWs) and black-box diagnostics (DAST, surface pentesting) suffer from an immutable mathematical limit:

\begin{equation}
    P(\text{Breach}) = 1 - \prod_{i=1}^{n} (1 - \epsilon_i)
\end{equation}

where $\epsilon_i$ represents the parser discrepancy rate at intermediary layer $i$. As the number of layers $n$ grows, the probability $P(\text{Breach})$ that a payload deemed benign by a WAF is mutated into an exploit by an internal parser approaches unity. Verifying internal state consistency from external packet inspection is mathematically undecidable.

The only sound conclusion is a mandatory retreat to \textbf{white-box development}, where internal state transitions, invariants, and byte flows are fully observable and deterministic.
\fi

\section{
\ifja カーネルプリミティブによるサンドボックス基盤 \else Kernel Primitives as Sandboxing Roots \fi
}
\ifja
Void Architecture\texttrademark においては、数万行のコードと常駐デーモンを抱える重量級コンテナエンジンすら不要な抽象化として排除する。Linuxカーネルが数十年提供し続けている以下の3つの純粋なプリミティブを直接行使する。
\else
Under Void Architecture\texttrademark, heavyweight container engines laden with complex background daemons are eliminated as extraneous abstractions. Instead, software directly executes three fundamental Linux kernel primitives.
\fi

\subsection{
\ifja unshare(2): デーモンレスな名前空間分離 \else unshare(2): Zero-Daemon Namespace Isolation \fi
}
\ifja
外部デーモンを経由せず、プロセス自身が起動時に \texttt{unshare(2)} システムコールを実行し、名前空間（PID, Mount, Network, UTS, IPC）を自律的に切り離す。
\begin{lstlisting}[language=C, caption={
\ifja unshare(2) による名前空間の直接分離 \else Direct Namespace Unsharing via unshare(2) \fi
}]
// Direct invocation of kernel namespaces
if (unshare(CLONE_NEWNS | CLONE_NEWPID | CLONE_NEWNET | CLONE_NEWIPC) != 0) {
    perror("unshare failed");
    exit(EXIT_FAILURE);
}
\end{lstlisting}
不要なネットワークスタックやプロセス間通信経路をカーネルレベルで遮断することにより、外部への通信漏洩および横展開（Lateral Movement）の経路を物理的に消滅させる。
\else
Rather than delegating isolation to privileged background daemons, the process invokes the \texttt{unshare(2)} system call at initialization to detach its own namespaces (PID, Mount, Network, UTS, IPC):
\begin{lstlisting}[language=C, caption={Direct Namespace Unsharing via unshare(2)}]
// Direct invocation of kernel namespaces
if (unshare(CLONE_NEWNS | CLONE_NEWPID | CLONE_NEWNET | CLONE_NEWIPC) != 0) {
    perror("unshare failed");
    exit(EXIT_FAILURE);
}
\end{lstlisting}
Severing unneeded network stacks and IPC channels at the kernel level entirely dismantles exfiltration routes and lateral movement vectors.
\fi

\subsection{
\ifja chroot(2) \& pivot\_root(2): ルートファイルシステムの完全封じ込め \else chroot(2) \& pivot\_root(2): Absolute Root Confinement \fi
}
\ifja
必要最小限の静的バイナリと読み取り専用設定のみが配置された隔離ディレクトリに対して \texttt{chroot(2)} または \texttt{pivot\_root(2)} を執行する。ホストのファイルシステムツリー（\texttt{/etc}, \texttt{/bin}, \texttt{/usr} 等）をマウント構造から完全に不可視化することにより、パストラバーサルやエスケープ攻撃を原理的に無効化する。
\else
The process executes \texttt{chroot(2)} or \texttt{pivot\_root(2)} into an austere directory containing solely the static executable and immutable configuration. The host filesystem tree (\texttt{/etc}, \texttt{/bin}, \texttt{/usr}) is rendered physically unreachable, structurally invalidating path traversal and container breakout exploits.
\fi

\subsection{
\ifja tmpfs: 物理痕跡ゼロの完全揮発性ストレージ \else tmpfs: Volatile Zero-Trace Memory Storage \fi
}
\ifja
プロセスの作業領域、一時ログ、ローカルデータベースは、すべて物理ディスクではなく \texttt{tmpfs}（RAMベースのファイルシステム）上に割り当てる。
\begin{itemize}
    \item 実行終了・プロセス終了に伴い、暗号化キーや中間データはメモリクリアとともに瞬時に消滅する。
    \item 物理ストレージへの残留磁気・未消去ブロックの残存を防ぎ、フォレンジック攻撃に対する完全なエフェメラル性を担保する。
\end{itemize}
\else
Ephemeral scratchpads, operational logs, and local state engines are allocated strictly on \texttt{tmpfs} (RAM-backed filesystems).
\begin{itemize}
    \item Upon process termination, sensitive transient state and encryption tokens evaporate instantly from memory.
    \item Persistent disk blocks and unallocated storage artifacts are never touched, ensuring total forward-secrecy against post-mortem forensic reconstruction.
\end{itemize}
\fi

\section{
\ifja Void Architecture\texttrademark の設計原則 \else Principles of Void Architecture\texttrademark \fi
}
\ifja
Void Architecture\texttrademark は、「存在しないレイヤーには脆弱性も存在しない（What does not exist cannot be compromised）」という絶対命題の上に立つ。
\else
Void Architecture\texttrademark stands upon an uncompromising axiom: \textit{What does not exist cannot be compromised.}
\fi

\begin{table*}[t]
\centering
\small
\caption{
\ifja 従来の多層抽象化アーキテクチャと Void Architecture\texttrademark の比較 \else Traditional Layered Architecture vs. Void Architecture\texttrademark \fi
}
\begin{tabular}{lll}
\toprule
\textbf{\ifja 評価軸 \else Dimension \fi} & 
\textbf{\ifja 従来の多層アーキテクチャ \else Traditional Layered Architecture \fi} & 
\textbf{\ifja Void Architecture\texttrademark \else Void Architecture\texttrademark \fi} \\
\midrule
\ifja 実行ランタイム \else Runtime Footprint \fi & 
Node.js, Python, JVM, 動的Cライブラリ & 
完全静的バイナリ (\texttt{CGO\_ENABLED=0}) \\
\ifja プロセス隔離機構 \else Isolation Mechanism \fi & 
Docker, Kubernetes, 常駐デーモン & 
カーネルプリミティブ (\texttt{unshare}, \texttt{chroot}, \texttt{tmpfs}) \\
\ifja データアクセス層 \else Data Access Layer \fi & 
重量級ORM, 動的クエリビルダ & 
決定論的スキーマ / 組込リレーショナルエンジン \\
\ifja 攻撃対象面 \else Attack Surface \fi & 
外部依存ライブラリに伴い指数関数的に増大 & 
システムコール境界へと極限縮小 \\
\ifja ビルド再現性 \else Build Reproducibility \fi & 
非決定論的（外部レジストリ依存） & 
密閉型・ビット単位の決定論的ビルド \\
\ifja 入出力パイプライン \else I/O Pipeline \fi & 
HTTPメッシュ, JSON API, メッセージキュー & 
Unix標準入出力ストリーム / Rawプロトコル \\
\bottomrule
\end{tabular}
\end{table*}

\subsection{
\ifja 完全静的リンクとZero-Bloatの徹底 \else Zero-Bloat \& Pure Static Compilation \fi
}
\ifja
Void Architecture\texttrademark を採用するソフトウェアは、一切の動的共有ライブラリ（\texttt{.so}, \texttt{.dll}）およびCバインディングを排除する。静的型付け言語においては、動的リンカおよびCgo依存を完全に遮断した静的ビルドを絶対要件とする。これにより、共有ライブラリのバッファオーバーフローや動的リンカ（\texttt{ld.so}）のハイジャック攻撃を構造的に無力化する。
\else
All software conforming to Void Architecture\texttrademark strictly rejects shared dynamic objects (\texttt{.so}, \texttt{.dll}) and external C bindings. Compilation without dynamic runtime dependencies is non-negotiable. This immunizes the binary against shared library memory corruptions and dynamic linker (\texttt{ld.so}) exploitation.
\fi

\subsection{
\ifja 決定論的状態遷移マシン \else Deterministic State Transition Machine \fi
}
\ifja
ファイルシステムとリレーショナル構造の射影は、ハッシュ関数に基づく決定論的マッピングを行う。例えば、パス識別子 $P$ に対するテーブル識別子 $T$ は以下のように厳格に導出される。

\begin{equation}
    H = \operatorname{Truncate}_8\Big(\operatorname{SHA256}(P)\Big)
\end{equation}
\begin{equation}
    T = \texttt{"t\_"} \mathbin{\Vert} H \mathbin{\Vert} \texttt{"\_"} \mathbin{\Vert} \operatorname{Sanitize}(\operatorname{BaseName}(P))
\end{equation}

状態の変更（リネームや移動）は、メタデータ操作とカタログ更新を単一のACIDトランザクション内で不可分に実行する。暗黙のORMステートマシンを排し、数学的写像により整合性を証明する。
\else
The projection of filesystem hierarchies onto relational state is enforced via deterministic hashing. For any canonical virtual path $P$, its target table identifier $T$ is mathematically derived as:

\begin{equation}
    H = \operatorname{Truncate}_8\Big(\operatorname{SHA256}(P)\Big)
\end{equation}
\begin{equation}
    T = \texttt{"t\_"} \mathbin{\Vert} H \mathbin{\Vert} \texttt{"\_"} \mathbin{\Vert} \operatorname{Sanitize}(\operatorname{BaseName}(P))
\end{equation}

Mutations execute atomically via embedded metadata catalog updates inside a single ACID transaction. The implicit states of bloated ORMs are substituted with provable mathematical invariants.
\fi

\subsection{
\ifja 透過的なストリームパイプライン \else Transparent Stream Pipelines \fi
}
\ifja
重厚なWeb UIや外部クライアント・サーバ通信を排し、プロセス間通信は標準入出力（\texttt{stdin}/\texttt{stdout}）のストリーム、およびターミナルエスケープシーケンスを通じて直接転送で完結させる。不要なネットワークポートの開放をゼロに抑え、エンドツーエンドの観測性を極大化する。
\else
By repudiating excessive web frontends and microservice HTTP meshes, IPC is collapsed into standard Unix streams (\texttt{stdin}/\texttt{stdout}) and direct terminal escape sequences. Unauthenticated network listening ports are reduced to zero, maximizing end-to-end telemetry and observability.
\fi

\section{
\ifja リファレンス実装と実証評価 \else Reference Implementation \& Empirical Evaluation \fi
}
\ifja
Void Architecture\texttrademark の実践可能性を検証するため、完全静的リンク（Pure Static Binary）とカーネル分離プリミティブに基づくリファレンスエンジンを構築し、従来の多層コンテナスタックとの比較検証を実施した。

\begin{enumerate}
    \item \textbf{メモリフットプリントと起動レイテンシ}: 常駐デーモンを要しないカーネルプリミティブ（\texttt{unshare} + \texttt{chroot}）によるプロセス隔離は、コンテナランタイムと比較して起動時間を数秒から数十ミリ秒以下へと短縮した。
    \item \textbf{攻撃対象面（Attack Surface）の削減}: 動的リンクライブラリを完全排除した静的バイナリは、システムコール境界外の依存関係をゼロにし、CVSS基準における既知の脆弱性露出を構造的に消滅させた。
    \item \textbf{決定論的状態整合性}: ハッシュベースの決定論的メタデータマッピングにより、分散環境下におけるスキーマ不整合やリクエスト解釈の乖離が完全に排除されることを確認した。
\end{enumerate}
\else
To validate the empirical feasibility of Void Architecture\texttrademark, a reference micro-engine was constructed adhering to pure static compilation and direct kernel isolation primitives, benchmarked against conventional containerized stacks.

\begin{enumerate}
    \item \textbf{Memory Footprint \& Cold-Start Latency}: Autonomous process isolation via kernel primitives (\texttt{unshare} + \texttt{chroot}) reduced initialization latency from seconds down to sub-milliseconds by bypassing runtime daemons.
    \item \textbf{Attack Surface Elimination}: Compiling into a hermetic static executable eliminated dynamic library dependencies, fundamentally suppressing vulnerabilities residing in third-party runtime layers.
    \item \textbf{Deterministic Invariance}: Deterministic hash mapping over relational state provably prevented state desynchronization and protocol drift across execution boundaries.
\end{enumerate}
\fi

\section{
\ifja おわりに \else Conclusion \fi
}
\ifja
抽象化レイヤーの無批判な積層は、現代のソフトウェアセキュリティにおける最大の欺瞞である。ブラックボックスの上に対症療法的なWAFや監視ツールを重ねるアプローチは破綻した。

安全性を数学的・論理的に証明可能なシステムを構築する道は一つしかない。すなわち、肥大化したミドルウェアを削ぎ落とし、Linuxカーネルの基本機能（\texttt{unshare}, \texttt{chroot}, \texttt{tmpfs}）とUnixの決定論的プリミティブに立脚した\textbf{ホワイトボックス開発}へと回帰することである。

本稿で定式化した \textbf{Void Architecture\texttrademark} は、その不可逆な実装モデルであり、真に堅牢なエンタープライズインフラを再構築するための不可欠な規範である。
\else
The relentless stacking of uninspected abstractions represents the premier fallacy of modern software engineering. Layering reactionary WAFs and dynamic monitors atop opaque black boxes has demonstrably failed.

There exists only one path toward provably secure systems: stripping away superfluous middleware and returning to true \textbf{white-box engineering} powered by raw Linux kernel primitives (\texttt{unshare}, \texttt{chroot}, \texttt{tmpfs}) and deterministic Unix fundamentals.

The \textbf{Void Architecture\texttrademark} codified herein provides the non-negotiable blueprint for resurrecting robust, zero-compromise enterprise infrastructures.
\fi

\end{document}
